\documentclass[11pt,a4paper]{article}
\usepackage[T1]{fontenc}
\usepackage[utf8]{inputenc}
\usepackage{amsmath,amssymb,xcolor}
\usepackage[margin=25mm]{geometry}
\usepackage[colorlinks=true,urlcolor=blue]{hyperref}
\definecolor{nexorablue}{HTML}{2459B5}
\setlength{\parindent}{0pt}
\setlength{\parskip}{7pt}
\setlength{\emergencystretch}{2em}
\renewcommand{\familydefault}{\sfdefault}
\begin{document}

{\small\bfseries\color{nexorablue} NEXORA CPGE}\par
{\small MPSI / NEX-MES-04 / Corrigé}\par
{\LARGE\bfseries Deux valeurs de la pesanteur sont-elles compatibles ?}\par
\vspace{0.25em}\hrule\vspace{0.6em}

\section*{Calcul}
L’incertitude de la différence vaut $\sqrt{0{,}07^2+0{,}04^2}\simeq0{,}0806\,\mathrm{m\,s^{-2}}$. Ainsi $z\simeq1{,}49$.

\section*{Conclusion}
Comme $1{,}49<2$, les valeurs sont compatibles pour le seuil annoncé. Le test ne prouve pas qu’elles sont exactement égales ni que les deux protocoles sont sans biais.


\end{document}
